\section{How many primes must a density rule couple?}
Let $\alpha$ be the supremum of natural densities of subsets of
$\mathbb N_{>0}$ containing no three distinct integers $a,b,c$ with
$ac=b^2$. Partition the primes into blocks of size at most $d$. For each
block choose an arbitrary set of allowed exponent vectors, and admit
an integer when its vector belongs to the allowed set in every block.
Let $\alpha_d$ be the corresponding supremum over all partitions and
rules giving a progression-free set with natural density.
The partition can vary with the competitor. In particular this model
can impose dependence among infinitely many separate groups of primes.

\medskip\noindent\textbf{Theorem (bounded prime blocks).}
For each $d\ge1$, the number $\alpha_d$ is computable and is attained.
An attaining construction uses Rankin's independent exponent rule
outside a finite set of primes. Allowing upper density gives the same
$\alpha_d$ in this class. Furthermore
\[
                  \alpha_d<\alpha,\qquad \lim_{d\to\infty}\alpha_d=\alpha .
\]
Thus every fixed bound on the number of jointly tested primes has a
strict density cost, even after optimizing the partition.
The result does not determine the unrestricted value $\alpha$ or the
first block size that improves on the one-prime rule.

\paragraph{Local probabilities.}
For a finite prime set $B$, the $B$-smooth positive integers form a
semigroup $S_B$. For $F\subset S_B$ write
\[
  c_B=\prod_{p\in B}(1-1/p),\qquad
  \mu_B(F)=c_B\sum_{s\in F}1/s.
\]
This is the natural density of integers with $B$-smooth part in $F$.
Indeed the finite exponent frequencies follow by the Chinese remainder
theorem, and discarding exponents $\ge m$ costs at most
$\sum_{p\in B}p^{-m}$ in upper density. Disjoint finite blocks have
independent probabilities. A nonempty progression-free product forces
each block rule to be progression-free: replace the smooth part of
one admitted integer by the three local terms to see this. Conversely
a product of progression-free rules is progression-free by projection.

Let $G$ consist of the nonnegative integers whose base-three digits
are zero or one, and let $C_B$ require all exponents in $G$. In an
arithmetic progression of three exponents in $G$, the lowest varying
digit would give a nonconstant progression in $\{0,1\}$ modulo three,
which is impossible. Thus $C_B$ is progression-free.
For infinite collections of Rankin primes the product of local
probabilities gives the natural density: the possible violations at
$p>Y$ require $p^2\mid n$ and cost at most $\sum_{p>Y}p^{-2}$.
The infinite product is positive since every factor is at least
$1-p^{-2}$.

\paragraph{Most blocks admit Rankin's rule as an optimizer.}
If $1/c_B\le2$, then $\mu_B(F)\le\mu_B(C_B)$ for every
progression-free $F$. This is the light-block observation following
McNew's Theorem~6.1. Here is a direct proof. For $s\in F\setminus C_B$
replace every ternary digit $2$ in its exponents by $0$, respectively
$1$, to get integers $a,b\in C_B$ with $a<b<s$ and $as=b^2$.
At least one of these proper divisors is missing from $F$. Charge
$1/s$ to one such divisor $t$. Each charged $s$ is a different
$t u$ with $u\in S_B\setminus\{1\}$. The charge at $t$ is at most
$(1/c_B-1)/t\le1/t$. Summing gives the claim.

Call the remaining blocks heavy. If every prime of $B$ exceeds $4d$,
then $\log(1/c_B)\le\sum_{p\in B}1/(p-1)\le1/4<\log 2$.
Hence at most $\pi(4d)$ blocks are heavy, and their union $H$
has at most $h=d\pi(4d)$ primes. Replace every light block by
Rankin's rule. The new set has natural density equal to the finite
heavy probability times the Rankin product off $H$. Its density is
at least the original upper density, since finite collections of
old block restrictions bound the latter from above, and the product
of the corresponding new probabilities converges to the new density.
A heavy rule may likewise be optimized separately. Its maximum exists:
indicators of a maximizing sequence on the countable semigroup admit
a pointwise convergent subsequence by diagonal extraction, avoidance
passes to the limit, and the weights $1/s$ are summable.

\paragraph{Existence of an optimal partition.}
Take a sequence approaching $\alpha_d$, already in the form just
described. Pass to a subsequence with constant number of exceptional
primes and constant partition of their ordered positions. Each of the
at most $h$ prime labels, after another subsequence, is either constant
or tends to infinity. In every block temporarily impose valuation zero
at each escaping prime $p$. The density loss is at most the sum of
$1/p$ over these primes, hence tends to zero. The surviving local
zero slices are progression-free on the fixed prime labels. One
may now use the optimal rule on each fixed block and Rankin's rule
at each discarded prime, since zero is an allowed Rankin exponent.
This single fixed construction attains at least the limiting density,
and belongs to the original class. The upper-density assertion follows
from the replacement bound above.

\paragraph{A uniform improvement.}
Fix the at most $h$ exceptional primes of a competitor after replacement.
A marker prime $\ell$ outside $H$ can be taken at most $L$, the
$(h+1)$st prime. Choose $Y\ge L$, depending only on $d$, so large that
\[
          \sum_{q\le Y,\ q\ {\rm prime}}\frac1q
                > 2L^2+(h+1)/2 .
\]
This is possible by Euler's divergence theorem. For the satellite set
$B$ use all primes through $Y$ outside $H\cup\{\ell\}$. Thus
$\sum_{q\in B}1/q>2L^2$. Apply the exchange of Section~\ref{sec:cap} to the lowest digits at the marker and satellites. Its cap proof works with any marker and leaves the higher digits binary.

Write $w=1/\ell$, $u_q=1/q$. For fixed good higher digits the change
in unnormalized reciprocal mass, relative to their common weight, is
\[
   \left(\prod_{q\in B}u_q\right)
       \left[w^2\left(1+\sum_{q\in B}u_q\right)-(1+w)\right]
      >\tfrac12\prod_{q\in B}u_q .
\]
The inequality uses $w\ge1/L$, $w\le1/2$. The higher digits have
total weight at least one. With $P_Y=\prod_{p\le Y}p$ and
$c_Y=\prod_{p\le Y}(1-1/p)$, the new local probability exceeds the
old one by at least $\delta_d=c_Y/(2P_Y)>0$. Since local probability
is at most one, this increases the full product by a factor at least
$1+\delta_d$. Thus $\alpha\ge(1+\delta_d)\alpha_d>\alpha_d$.
In particular, even allowing each partition to use distant, arbitrarily
large primes gives a uniform gap for fixed $d$.

\paragraph{Approximation by a finite block.}
Let $P$ be the primes through $Y\ge2$, let $c=c_P$ and let $M$ be the
largest reciprocal mass $\sum_{s\in F}1/s$ for a progression-free
$F\subset S_P$. For any progression-free set $A$ of natural density,
write $n=rs$ with $(r,\prod_{p\in P}p)=1$ and $s\in S_P$.
For each fixed $r$ the permitted $s$ form a progression-free set.
Consequently
\[
 \sum_{n\le X,\,n\in A}\frac1n
  \le M\sum_{r\le X,\,(r,\prod_{p\in P}p)=1}\frac1r
  =M(c\log X+O_Y(1)).
\]
The last estimate follows by finite inclusion--exclusion; partial
summation gives $d(A)$ as the limit of the left side divided by
$\log X$. This proves $\alpha\le cM$. A maximizing smooth template
on $P$ and Rankin's rule elsewhere, which has width $\pi(Y)$, gives
\[
  \alpha_{\pi(Y)}\ge cM\prod_{p>Y}(1-p^{-2})
       \ge cM\,\frac{Y}{Y+1}
\]
for integer $Y$, by telescoping over all integers greater than $Y$.
Together with monotonicity this proves convergence. Choosing $Y=p_{d+1}-1$, where $p_j$ is the $j$th prime, gives $0<\alpha-\alpha_d\le 1/p_{d+1}$ since $cM\le1$. Thus $\alpha$ itself is computable.

\paragraph{Effective computation for fixed width.}
Uniformly over competitors after replacement, freezing exponents at
heavy primes $p>Y$ to zero costs at most $h/Y$. Releasing just
these primes to use Rankin's rule can only increase density. Thus
optimizing partitions of primes through $Y$ of width $d$, with the
Rankin tail, approximates $\alpha_d$ from below with error at most
$h/Y$. There are finitely many partitions to search. Truncate each
prime exponent below $m$; in total this costs at most
$\sum_{p\le Y}p^{-m}$. On finite boxes, enumerate all subsets,
reject any nonconstant arithmetic progression of exponent vectors,
and maximize the exact rational probability. Finally the Rankin
tail product through any $Z>Y$ has an absolute truncation error at most
$\sum_{n>Z}n^{-2}\le1/Z$; its single-prime factors can be truncated
geometrically. These estimates give rational intervals of arbitrary
prescribed width. They make no assertion of practical running time.

\paragraph{Relation to earlier work.}
McNew's Theorems 5.1--5.2 concern unrestricted \emph{upper} density,
and his Theorem~6.1 treats a one-prime logarithmic optimum.
See \href{https://www.nathanmcnew.com/GPFsets.pdf}{McNew's paper,
Sections 5--7}. The incoming manuscript
\texttt{mcnew\_counterexample\_and\_grh\_transfers.zip} disproves
Conjecture~7.2 by a marker and a shared deletion and shows that any
fixed finite change to Rankin's rule admits improvement. The bounded-block
theorem allows an arbitrary rule in every one of infinitely many
blocks, proves a gap uniform over all their partitions, and gives
attainment and effective approximation for each width. The lowest-digit
construction above supplies its local gain.
